% ch01-question — the central question and the map from questions to models

\section{Worth, formally}\label{ch01:sec:worth}

The business objective is to understand which private-market investments are
``worth it.'' That phrase hides a decision problem, and the first discipline
this document imposes is to make the decision problem explicit
\emph{before} any model is chosen. An investment thesis about a private
company is a bet on a small set of future observables: whether the company
raises again, when it raises, at what valuation, and whether the eventual
exit repays the capital deployed today. A model earns its keep only by
improving one of those bets. So the central question is not ``which model is
best'' but:

\begin{center}
\emph{What is valuable to know?}
\end{center}

Formally, let $\mathcal{D}$ be a space of actions available to the investor
(allocate attention to a company, enter a diligence process, price a
secondary position, pass), let $\omega$ denote the future evolution of the
market --- the marked point process of deals defined in \cref{ch:data} ---
and let $u(d,\omega)$ be the utility of action $d\in\mathcal{D}$ under
outcome $\omega$. A predictive object $\pi$ (a probability, an intensity, a
ranking) is \emph{valuable} exactly to the extent that
\begin{equation}\label{eq:ch01:voi}
\mathrm{VoI}(\pi)
 \;=\;
 \E\!\left[\, u\!\left(d^{*}_{\pi},\,\omega\right) \right]
 \;-\;
 \E\!\left[\, u\!\left(d^{*}_{0},\,\omega\right) \right],
 \qquad
 d^{*}_{\pi} = \argmax_{d\in\mathcal{D}} \E\!\left[u(d,\omega)\st \pi\right],
\end{equation}
is positive and large, where $d^{*}_{0}$ is the action taken without the
model. This is the standard value-of-information decomposition, and its
practical force here is negative: it rules out building models whose outputs
no action consumes. Every model family developed in this document is
admitted only because a concrete action consumes its output, and the
evaluation metric of each family is chosen to be a proxy for
\cref{eq:ch01:voi} under that action, not a generic accuracy score.

\section{The four answers and the targets they induce}\label{ch01:sec:targets}

Interrogating the sourcing and monitoring workflows that would consume model
output yields four distinct answers to ``what is valuable to know,'' each
inducing a different mathematical target on the same underlying data.

\paragraph{T1 --- Identity: \emph{which} company gets funded next.}
The sourcing action is a short list: given finite diligence capacity, which
$K$ companies should be looked at now because they are most likely to be the
next ones to raise? The natural target is the conditional distribution of
the next deal's identity,
\begin{equation}\label{eq:ch01:t1}
\Prob\big(\text{next deal in sector } s \text{ goes to firm } i
 \;\big|\; \Ftm\big),
\end{equation}
and the natural metric is ranking quality --- recall at $K$ against the next
deals actually observed (\cref{ch:ranking}). Crucially, \cref{eq:ch01:t1}
conditions on \emph{a deal occurring}: the timing of deals cancels from the
target, which is why identity models can survive data defects that destroy
timing models (\cref{ch01:sec:factorization}).

\paragraph{T2 --- Timing: \emph{when} the deals arrive.}
Pricing, portfolio pacing, and capital-deployment planning consume arrival
times: how many deals will sector $s$ produce next quarter, and how is firm
$i$'s next round distributed in time? The target is the conditional
intensity $\lam(t)$ of the deal process (\cref{ch:foundations}) and
functionals of it (expected counts, waiting-time distributions), evaluated
by likelihood on held-out windows and by calibration diagnostics
(\cref{ch:gof}). This is the target the Poisson and Hawkes machinery of
\cref{ch:poisson,ch:hawkes} serves.

\paragraph{T3 --- Horizon events: will firm $i$ raise within $\horizon$.}
Many consuming actions are binary and horizon-bound: enter diligence now if
the company is likely to raise within the next six months; deprioritize if
no event is plausible within two years. The target is the fixed-horizon
event probability
\begin{equation}\label{eq:ch01:t3}
p_{i}(t,\horizon)
 \;=\;
 \Prob\big(N_{i}(t+\horizon) - N_{i}(t) \ge 1 \st \Ft\big),
 \qquad \horizon \in \{\text{6M},\, \text{1Y},\, \text{2Y}\},
\end{equation}
a discrete-time object that can be attacked directly with supervised
learning on a bucketed panel (\cref{ch:buckets}) --- at the price of label
construction subtleties that \cref{ch:buckets} treats as first-class
mathematics, not plumbing.

\paragraph{T4 --- Marks: size, valuation, terms.}
Finally, worth depends on \emph{what} the deal is, not only that it happens:
round size, pre/post valuation, up/down/flat classification. These are mark
distributions conditional on an event, $\Prob(\dmark \in \cdot \st
\text{deal at } t \text{ by } i, \Ftm)$. This document develops T4 only to
the extent needed by the decision layer (\cref{ch:decision}); its
estimation reuses the regression machinery of \cref{ch:buckets} and
inherits every covariate caveat of \cref{ch:stale}.

These four targets are not competing models of the same quantity; they are
different marginals and conditionals of one object --- the marked point
process --- consumed by different actions. The recurring failure mode this
document is designed to prevent is answering one target's question with
another target's model and metric (for instance, selling a horizon
classifier's probability, T3, as a timing model, T2; or judging a ranking
model, T1, by likelihood alone).

\section{One process, two factorizations}\label{ch01:sec:factorization}

The unifying object is the marked point process of deals: ground times $t_m$
(when a deal happens), marks $(s_m, i_m, \dmark_m)$ (which sector, which
firm, what terms). Its likelihood factorizes \emph{exactly} as
\begin{equation}\label{eq:ch01:groundmark}
p\big(\text{times, sectors, firms, terms}\big)
 \;=\;
 \underbrace{p\big(\text{times, sectors}\big)}_{\text{ground process}}
 \;\times\;
 \underbrace{p\big(\text{firm} \st \text{sector, time, history}\big)}_{\text{mark factor}}
 \;\times\;
 \underbrace{p\big(\text{terms} \st \text{firm, time}\big)}_{\text{deal marks}} .
\end{equation}
This identity --- established as \cref{thm:ch03:groundmark} and inherited
from the repository's earlier theory notes, where it is the decomposition
labeled M1 --- is the load-bearing wall of the whole design. It is an
identity of likelihoods, not an approximation: all approximations live
\emph{inside} one factor at a time, and defects of the data can be assigned
to the factor they damage. Three consequences organize the document.

First, \emph{the targets map onto factors}. T2 lives in the ground process;
T1 lives in the mark factor; T4 lives in the deal marks; T3 is a horizon
functional of ground $\times$ mark. A team can therefore improve one factor
without touching the others, and \cref{ch:roadmap} assigns code modules
factor by factor.

Second, \emph{the mark factor is robust where the ground process is
fragile}. The mark factor is a conditional-choice likelihood in which the
sector-time baseline cancels (\cref{ch:ranking}); it survives untrusted
deal dates, weekly aggregation, and even moderate intensity misspecification.
The ground process absorbs precisely those defects. When the data cannot
support a trustworthy ground model, the mark factor still delivers T1 ---
this asymmetry, proved and quantified in
\cref{ch:censoring,ch:ranking}, is why ranking is the most defensible
deliverable on current data quality.

Third, \emph{discrete-time models are the same object at coarser
resolution}. The bucketed panel of \cref{ch:buckets} is the ground+mark
object aggregated to intervals: \cref{thm:ch09:hazardlink} shows the
bucket-probability model is exactly a discrete-time hazard model, and
recovers the continuous-time intensity as bucket width shrinks. XGBoost on
buckets and Hawkes in continuous time are therefore rungs of one ladder,
not rival ideologies; which rung to stand on is decided by the target and
by what the data defects of \cref{ch:data} leave standing.

\section{From questions to models: the decision map}\label{ch01:sec:map}

\Cref{tab:ch01:map} is the document in one table. Each row is an answer to
``what is valuable to know''; the columns give the mathematical target, the
model family this document commits to for that target, the evaluation
metric that proxies \cref{eq:ch01:voi}, and the chapter that develops it.

\begin{table}[htbp]
\centering
\small
\begin{tabular}{p{0.21\textwidth} p{0.2\textwidth} p{0.23\textwidth} p{0.16\textwidth} p{0.07\textwidth}}
\toprule
\textbf{Valuable to know} & \textbf{Target} & \textbf{Committed model} &
\textbf{Metric} & \textbf{Ch.} \\
\midrule
Which firm raises next (identity) &
$\Prob(i \mid \text{deal}, s, t, \Hist)$ &
risk-set conditional logit $\to$ gradient-boosted ranker &
recall@$K$, MRR, held-out conditional NLL & \ref{ch:ranking} \\
\addlinespace
When/how many deals (timing) &
ground intensity $\lam_s(t)$ &
covariate NHPP $\to$ exponential Hawkes with engineered
excitation/inhibition &
held-out NLL, time-rescaling, PIT & \ref{ch:poisson}--\ref{ch:gof} \\
\addlinespace
Will $i$ raise within $\horizon$ (horizon) &
$p_i(t,\horizon)$, $\horizon\in\{$6M, 1Y, 2Y$\}$ &
discrete-hazard panel + XGBoost with staleness features &
calibrated probabilities; PR-AUC; lift@$K$ & \ref{ch:buckets} \\
\addlinespace
Deal terms (marks) &
$p(\dmark \mid i, t)$ &
mark regressions on deal covariates &
per-mark scores & \ref{ch:decision} \\
\bottomrule
\end{tabular}
\caption{The decision map: answers to ``what is valuable to know,'' the
targets they induce, and the models and metrics this document commits to.}
\label{tab:ch01:map}
\end{table}

\begin{decision}{D1.1 --- target before model}
No model is built, modified, or benchmarked in this repository without first
naming which target (T1--T4) it serves and which action consumes its output.
Pull requests that improve a metric not attached to a target are rejected by
policy. Rationale: \cref{eq:ch01:voi}; the failure mode it prevents is
optimizing likelihood in a factor whose output nothing consumes.
Reversal trigger: none anticipated --- this is the document's charter.
\end{decision}

\begin{decision}{D1.2 --- both stacks, one object}
We maintain \emph{both} the continuous-time stack (NHPP/Hawkes ground
process + risk-set mark factor) and the discrete-time stack (bucketed
panel + XGBoost), as views of the same marked process at different
resolutions, rather than choosing one ideology. Rationale: the synthetic
rehearsals recorded in \cref{app:results} show the covariate-driven
baseline carries most of the predictive mass (nats-scale gains) while
excitation contributes decimal points --- but also show a parity XGBoost
beating the linear conditional logit on identity (the fired B1 test). Each
stack dominates somewhere. Reversal trigger: if, after the stale-covariate
remedies of \cref{ch:stale}, one stack dominates on \emph{every} consuming
action for two consecutive evaluation campaigns, the other is demoted to
benchmark-only status.
\end{decision}

\section{What ``worth it'' will mean operationally}\label{ch01:sec:operational}

The decision layer of \cref{ch:decision} composes the four targets into the
quantities the investment workflow actually consumes: a ranked sourcing
list with horizon-qualified event probabilities attached (T1 + T3), sector
pacing forecasts with uncertainty bands (T2), and expected-terms context
(T4). Until that chapter, ``worth it'' is used in the narrow, precise
sense of \cref{eq:ch01:voi}: a model output is worth producing if and only
if the action that consumes it gains utility, and the document's evaluation
protocols are designed so that offline metrics are monotone proxies for
that gain. The chapters in between are the machinery that makes those
outputs trustworthy on the data we actually have --- which is the subject
of the next chapter, and the source of every hard problem in this
document.

\begin{codehooks}
The decision map of \cref{tab:ch01:map} is mirrored in code by keeping the
factors of \cref{eq:ch01:groundmark} in separate modules with separate
evaluation harnesses: ground-process fitters (\modrepo{sector\_stability.py},
\modrepo{eventtime/}), mark-factor fitters (\modrepo{sector\_ranker.py},
\modrepo{models/hawkes\_ranker.py}), and (to be written) the bucketed-panel
builder and horizon classifier of \cref{ch:buckets}
(\modrepo{panel/}), the staleness state layer of \cref{ch:stale}
(\modrepo{covstate/}), and the decision layer of \cref{ch:decision}
(\modrepo{decision/}). \Cref{ch:roadmap} gives the full mapping with
migration steps.
\end{codehooks}
